\documentclass{../latex/utmreport}
\newcommand{\ReportTitle}{Lab 1: Concurrency in Practice}
% Shared authorship and metadata for this and all subsequent PR lab reports.
\newcommand{\StudentName}{Gafenco Victor and Balan Gabriel}
\newcommand{\StudentGroup}{FAF-241}
\newcommand{\CourseName}{Network Programming (PR)}
\newcommand{\InstructorName}{Fiștic Cristofor}
\newcommand{\InstructorTitle}{Assistant lecturer}
\newcommand{\ReportCity}{Chișinău}
\newcommand{\ReportYear}{2026}
\newcommand{\CourseRepository}{https://github.com/yoda-dinmd/pr-labs}

\selectlanguage{english}
\utmlogo{../latex/assets/UTM_Logo_English.jpeg}
\institution{Technical University of Moldova}
\faculty{Faculty of Computers, Informatics, and Microelectronics}
\department{Department of Software Engineering}
\studyprogramme{Software Engineering Study Programme}
\documenttype{Laboratory report}
\student{\StudentName}{\StudentGroup}
\renewcommand{\UTMStudentLabel}{Students}
\supervisor{\InstructorName}{\InstructorTitle}
\renewcommand{\UTMCoordinatorLabel}{Instructor}
\cityyear{\ReportCity}{\ReportYear}
\reportsubtitle{\CourseName}
\hypersetup{pdfauthor={\StudentName},pdfsubject={\CourseName},
  pdftitle={\ReportTitle}}
\reporttitle{\ReportTitle}
\addto\captionsenglish{\renewcommand{\figurename}{Figure}\renewcommand{\tablename}{Table}\renewcommand{\bibname}{References}}
\renewcommand{\figurename}{Figure}
\renewcommand{\tablename}{Table}
\renewcommand{\bibname}{References}
\usepackage{float}
\usepackage{placeins}
\setlength{\emergencystretch}{2em}
\newcommand{\code}[1]{\texttt{\detokenize{#1}}}
% Keep screenshots beside their observations, with bounded height and width.
\newcommand{\screenshot}[4][\textwidth]{%
  \begin{figure}[H]
    \centering
    \includegraphics[width=#1,height=0.62\textheight,keepaspectratio]{screenshots/#2}
    \caption{#3}\label{#4}
  \end{figure}}

\newcommand{\resultplot}[3]{%
  \begin{figure}[H]\centering
  \includegraphics[width=\textwidth]{results/plots/#1.png}
  \caption{#2}\label{#3}
  \end{figure}}

\begin{document}
\makeutmreporttitlepage
\tableofcontents
\clearpage
\begin{utmabstract}[Abstract]
This laboratory investigates how work distribution and synchronization affect
execution time, correctness, and resident memory when counting the signs of
50 million integers. We examine independent file partitions, three shared-array
counting strategies, a bounded producer--consumer queue, and eight files read
from USB storage. The saved measurements show a best speed-up of 2.857 for
independent partitions and 1.092 for the single-reader queue implementation. Local counters
complete in 269.903 ms, while a mutex around every shared update takes
11,213.016 ms. All three unsynchronized runs lose more than 35 million updates.
The USB series improves from 11.174 s to 6.522 s with four threads.
A separate six-reader experiment is retained for comparison: it has lower
elapsed times, although the runs were collected in different sessions. Conclusions
are restricted to the existing data; no new benchmark was executed for this
report.
\utmkeywords{concurrency, threads, synchronization, race condition, bounded queue, speed-up}
\end{utmabstract}

\chapter*{Abbreviations}
\addcontentsline{toc}{chapter}{Abbreviations}
\noindent\begin{tabularx}{\textwidth}{lY}
CPU & Central processing unit \\
I/O & Input/output \\
RSS & Resident set size \\
SMT & Simultaneous multithreading \\
USB & Universal Serial Bus \\
MiB & $2^{20}$ bytes \\
\end{tabularx}

\utmintroduction
Concurrency allows several activities to progress during the same interval;
parallelism executes work simultaneously on available processor resources.
Adding threads can reduce elapsed time, but shared locks, storage bandwidth,
allocation, and scheduling can limit the benefit. Our experiment uses one
C++ program with command-line selections for each counting strategy. The
work was completed jointly by Gafenco Victor and Balan Gabriel, group FAF-241.

We analyse the existing source, reference counts, CSV measurements, and supplied
plots. The objective is to identify where parallelism helps, quantify its
memory cost, and explain why correct synchronization and fast execution must
be considered together. Numerical observations are distinguished from possible
causes that would require profiling to establish.

\chapter{Environment and measurement method}
\section{Hardware and software}
Table~\ref{tab:environment} records the available local environment and build
metadata. This inventory was inspected when preparing the report; the CSV files
do not embed a historical hardware or compiler inventory. It therefore describes
the available workstation, without proving every setting at measurement time.
\begin{table}[H]\centering
\caption{Available environment and evidence}\label{tab:environment}
\begin{tabularx}{\textwidth}{|p{0.29\textwidth}|Y|}\hline
Item & Recorded value \\\hline
Processor & Intel Core i7-8650U, nominal 1.90 GHz \\\hline
Execution resources & 4 physical cores, 8 logical processors; 2 SMT threads per core \\\hline
Memory & Approximately 15 GiB usable RAM reported by the operating system \\\hline
Local storage & KINGSTON SNV2S1000G, NVMe SSD, 931.5 GiB \\\hline
Operating system & openSUSE Tumbleweed, snapshot 20260901; Linux 7.2.7-1-default, x86-64 \\\hline
Language and compiler & C++23; cached compiler: Clang 23.1.2 \\\hline
C++ runtime & Installed libstdc++6 package: 16.2.1+git9713-1.1 \\\hline
Build configuration & Cache says Release; saved compilation flags contain no optimization flag \\\hline
USB device and connection & Model, USB version, and negotiated link speed are not recorded in the supplied results \\\hline
\end{tabularx}
\end{table}
The saved flags include \code{-std=c++23} and warning options but no
\code{-O2} or \code{-O3}. We cannot establish the measured binary's optimization
level from a Release cache entry alone. The results should not be presented as
verified optimized-build performance. Task 1b records eight threads, consistent
with the workstation's logical processor count.

\section{Input and correctness}
The input file contains 50,000,000 integers in the range $[-1,000,000,1,000,000]$,
one per line. The generator uses \code{std::mt19937} with seed 42. The supplied
file occupies 369,436,719 bytes (approximately 352.32 MiB).
The reference and all correct-version CSV rows agree on the following counts:
\begin{center}
\begin{tabular}{lr}
Negative & 24,994,850 \\
Zero & 26 \\
Positive & 25,005,124 \\\hline
Total & 50,000,000 \\
\end{tabular}
\end{center}
Comparing all three counters is stronger than checking only the sum: equal
and opposite classification errors could preserve the total. The unsynchronized
version is intentionally excluded from the correctness requirement. Task 2
uses eight consecutive files with 6,250,000 integers each.

\section{Timing, memory, and aggregation}
The benchmark launcher starts a fresh counting process for each measurement.
The program uses \code{steady_clock}; timing includes task setup, thread
creation, counting, joining, and reduction. For Task 1b, loading the integer
array occurs before timing. Peak memory is the process-lifetime peak RSS,
so it still includes that array. Launcher and plotting memory are excluded.

Tasks 1a and 1c have three runs per configuration. Before each run, the launcher
reads the file without parsing to warm the cache. Their results therefore do
not measure cold local-disk throughput. Task 1b has one local-counter run,
one mutex run, and three race runs. Task 2 has one run at each thread count;
its CSV policy is \code{manual-reconnect}. The launcher requests ejecting,
unplugging, and reconnecting the USB before each run, without a warm-up read.
The CSV records that policy but does not prove the physical procedure occurred.

Tables report median elapsed time and median peak RSS where repeats exist.
Speed-up is computed from median times:
\begin{equation}
S(p)=\frac{\operatorname{median}(T(1))}{\operatorname{median}(T(p))}.
\end{equation}
For the main Task 1c series, $p$ means consumers, with one additional reader.
The archived experimental series uses six readers and is identified separately. The charts show individual samples and observed minimum--maximum
ranges, not confidence intervals. Task 1b compares strategies at a fixed eight
threads; no single-thread Task 1b baseline exists, so a thread-scaling speed-up
cannot be calculated for it.

\chapter{Task 1: Counting one dataset}
\section{Task 1a: Independent file partitions}
The coordinator obtains the file size and chooses approximately equal byte
ranges. Each internal boundary is advanced through the next newline. The same
boundary is the end of the preceding range and the start of the next range.
Workers open their own file streams, seek to their starts, parse complete lines,
and update local counters. The coordinator joins the workers and sums their
results. It does not load or parse the entire file before distributing work.

The ranges are contiguous and half-open: $[b_i,b_{i+1})$. A worker reads a line
only when its starting position precedes its end boundary. Thus the boundary
line belongs to exactly one range. This reasoning assumes the generated
newline-terminated input, for which byte positions advance by the line length
plus one. Equal bytes need not mean exactly equal numbers of lines or equal
execution time.

\resultplot{task-1a}{Task 1a: elapsed time, speed-up, and peak resident memory}{fig:1a}
\begin{table}[H]\centering
\caption{Task 1a medians from three runs per configuration}
\begin{tabular}{rrrr}\hline
Threads & Time (s) & Speed-up & Peak RSS (MiB) \\\hline
1 & 9.942 & 1.000 & 4.20 \\
2 & 5.995 & 1.658 & 4.20 \\
4 & 4.423 & 2.248 & 4.27 \\
8 & 3.971 & 2.504 & 4.08 \\
16 & 3.932 & 2.529 & 4.14 \\
32 & 3.767 & 2.639 & 4.08 \\
64 & 3.480 & 2.857 & 4.83 \\\hline
\end{tabular}
\end{table}

\subsection{Scaling and processor resources}
The largest improvements occur up to four threads, corresponding to the four
physical cores. Moving from four to eight threads reduces median time by about
10.2\%, much less than doubling throughput. Logical processors share a core's
execution resources and do not provide the capacity of extra physical cores.

There is no strict stopping point in the measured range: the best median occurs
at 64 threads. The curve flattens substantially around 8--16 threads, but
32 and 64 threads still improve these measurements. Relative to eight threads,
64 threads reduce median time by about 12.4\%. We therefore cannot claim that
32 and 64 threads were slower. Extra threads cannot create extra cores, and
context switching, contention, and cache pressure can eventually outweigh their
benefit. The modest continuing improvement could involve scheduling or load
balance; timings alone do not identify its cause. The eight-thread samples
span 3.799--4.875 s, illustrating considerable run variability.

\subsection{Memory and workload classification}
Median RSS stays near 4.1--4.3 MiB through 32 threads and rises to 4.83 MiB at
64. The program streams its partitions instead of retaining all input.
Additional thread stacks and stream buffers can increase memory, but reserved
stack address space is not equivalent to resident memory. Small RSS differences
also reflect allocation and page residency.

Warm-cache input and the substantial benefit of multiple cores are consistent
with a CPU-side parsing workload rather than slow physical-disk reads being
the dominant cost. This is an inference from the procedure and scaling, not a
CPU-utilization or storage-throughput profile.

\section{Task 1b: Shared and local counters}
All variants scan the same preloaded integer array using eight threads.
Version 1 updates three shared counters without synchronization. Version 2
locks a common mutex for every update. Version 3 keeps counters local to each
worker and reduces them after joining.
\resultplot{task-1b}{Task 1b: strategy timing, peak memory, and correctness}{fig:1b}
\begin{table}[H]\centering
\caption{Task 1b measurements; all configurations use eight threads}
\begin{tabular}{lrrr}\hline
Strategy / run & Time (ms) & Peak RSS (MiB) & Counted total \\\hline
Local counters & 269.903 & 194.82 & 50,000,000 \\
Mutex & 11,213.016 & 194.65 & 50,000,000 \\
Race, run 1 & 742.858 & 194.96 & 14,801,194 \\
Race, run 2 & 827.543 & 194.68 & 14,601,519 \\
Race, run 3 & 908.762 & 194.85 & 14,560,480 \\\hline
\end{tabular}
\end{table}
The measured ranking is local counters, unsynchronized counters, then mutex.
The median race time is 827.543 ms, about 3.07 times the local-counter time.
The mutex run takes about 41.54 times as long as the local-counter run.
Only the mutex and local variants match all three reference counts.
The race runs retain only about 29.1--29.6\% of the expected total.

An increment is conceptually a read, an addition, and a write. If two threads
read a counter at 10, both can compute 11 and both store 11; two updates have
increased the counter by only one. This illustrates lost updates, but the C++
data race has undefined behavior, so the language guarantees neither this
particular execution nor a particular wrong total. Even a matching total would
not establish correctness.

The mutex preserves updates but serializes 50 million critical sections and
requires repeated lock acquisition. Local counters remove this shared hot spot;
only the final reduction combines results. Unsynchronized shared writes can
still incur cache-coherence traffic, explaining why avoiding a lock need not
outperform local accumulation. The integer array alone occupies about
190.73 MiB when integers are four bytes; this accounts for most of the roughly
195 MiB peak RSS in every variant.

\section{Task 1c: Bounded queue and consumers}
The main experiment uses one reader and a bounded queue of at most 100 batches,
with 10,000 raw lines per batch. The reader builds batches of strings; consumers
remove batches, parse their lines outside the queue lock, and count locally.
A mutex protects the queue and the remaining-reader counter. Condition variables
block the reader when the queue is full and consumers when it is empty.
Consumers exit only when the queue is empty and the reader has finished.
The source now sets \code{READER_COUNT = 1}; the new CSV records total background
threads as workers plus one, excluding the main coordination thread.

\resultplot{task-1c}{Task 1c: one reader, variable consumers, and a bounded queue}{fig:1c}
\begin{table}[H]\centering
\caption{Single-reader Task 1c medians from three runs per configuration}
\begin{tabular}{rrrr}\hline
Workers & Time (s) & Speed-up & Peak RSS (MiB) \\\hline
1 & 15.781 & 1.000 & 5.09 \\
2 & 14.496 & 1.089 & 5.16 \\
4 & 14.518 & 1.087 & 5.10 \\
8 & 14.499 & 1.088 & 4.99 \\
16 & 14.561 & 1.084 & 5.15 \\
32 & 14.448 & 1.092 & 5.33 \\
64 & 15.416 & 1.024 & 5.57 \\
\hline
\end{tabular}
\end{table}

\subsection{Why adding consumers has little effect}
One consumer takes a median 15.781 s. Two consumers reduce this to 14.496 s,
a speed-up of only 1.089. From two through 32 consumers, medians remain between
14.448 and 14.561 s. The lowest median is at 32 consumers, but its difference
from two consumers is only about 0.34\%; this does not establish a meaningful
advantage. The practical plateau begins around two consumers. At 64 consumers,
the median rises to 15.416 s. The 32-consumer samples span 14.385--18.678 s,
showing why the minimum median should not be treated as a precisely located
optimum.

The reader still performs 50 million line reads and constructs the batches of
strings sequentially. Consumers only parse those short integer strings and
increment local counters. Once consumers can keep up with production, adding
more of them cannot increase the reader's rate. This is consistent with a
producer bottleneck; it is not evidence that the laptop cannot run threads
in parallel. Warm-cache input reduces physical storage reads but does not
remove line handling, string operations, queue transfers, or scheduling costs.
The inspected compilation flags also contain no optimization option, which is
another plausible contributor to absolute execution time. We did not change
build settings or collect a profiler trace, so individual costs are not isolated.

Median RSS remains between about 4.99 and 5.57 MiB, much lower than in the
six-reader experiment. This is consistent with a queue that does not retain
many batches because consumers keep up with the producer. Queue occupancy
and thread utilization were not measured, so this remains an interpretation
rather than a directly observed mechanism.

\subsection{Comparison with independent partitions}
Task 1a with one thread takes 9.942 s, compared with 15.781 s for Task 1c with
one consumer and one reader. At four workers/threads, Task 1a takes 4.423 s
versus 14.518 s for Task 1c. Their best recorded medians are 3.480 s and
14.448 s, respectively, a ratio of approximately 4.15. The runs were obtained
in different sessions, so these are descriptive comparisons, not controlled
estimates of the cost of changing only the algorithm.

Task 1a distributes both reading and parsing among workers. Task 1c leaves
line reading and batch construction with one producer and adds string storage
and queue coordination. This difference explains why extra consumers may have
little work to do even when there are unused processor resources.

\subsection{Additional experiment: six readers}
We also tried six readers to investigate whether parallel production changes
throughput. Each reader processes a disjoint, line-aligned file range and sends
batches to the same bounded queue. The archived CSV records workers plus six
background threads. This was an intentional extension of the main experiment.
Its plot was generated with labels that still say ``one reader''; for this
archived figure only, the actual reader count is six.
\resultplot{task-1c-6-readers}{Archived six-reader experiment; the image's one-reader labels are stale}{fig:1c-six}
\begin{table}[H]\centering
\caption{Recorded median times with one and six readers, at equal consumer counts}
\begin{tabular}{rrrr}\hline
Workers & One reader (s) & Six readers (s) & Time ratio \\\hline
1 & 15.781 & 7.614 & 2.07 \\
2 & 14.496 & 5.602 & 2.59 \\
4 & 14.518 & 5.520 & 2.63 \\
8 & 14.499 & 5.473 & 2.65 \\
16 & 14.561 & 5.206 & 2.80 \\
32 & 14.448 & 5.228 & 2.76 \\
64 & 15.416 & 5.213 & 2.96 \\
\hline
\end{tabular}
\end{table}

These saved results do not show that adding readers makes no difference.
At two consumers, the recorded median falls from 14.496 s with one reader to
5.602 s with six, a ratio of about 2.59. At 16 consumers, it falls from
14.561 s to 5.206 s, a ratio of about 2.80. The six-reader series itself gains
little from adding consumers beyond two, with its best median at 16 consumers
and essentially no further improvement at 32 or 64.

The distinction is between increasing the number of readers and increasing
the number of consumers. The recorded six-reader times are substantially lower,
while extra consumers have diminishing returns in both series. This is consistent
with production limiting the single-reader design. However, the series were
collected in separate sessions without a preserved, matched record of build
flags, system load, temperature, or CPU frequency. Consequently the ratios do
not establish the isolated causal benefit of adding five readers.

The archived peak RSS is about 47 MiB at one or two consumers and about 12 MiB
at 16--64 consumers. Multiple producers can retain more text in the queue or
in batches under construction. This agrees with increased production capacity,
although queue occupancy was not logged.

\subsection{Batch size and queue capacity}
The 100-batch queue limit bounds queued work to one million lines, but batches
under construction or being processed occupy additional memory. Single-line
batches would require up to 50 million queue transfers instead of roughly
5,000 full batches, increasing locking, notification, and allocation overhead.
An unbounded queue would remove backpressure: if production outpaces consumption,
much of the input could accumulate as strings and potentially cause paging.
A bounded queue controls that risk; its capacity is a maximum, not a promise
that the queue will fill during a run.

\chapter{Task 2: Files on USB storage}
Each thread owns a fixed selection of the eight files. With one, two, or four
threads, each processes eight, four, or two files respectively. A file is read
using one stream read request into one byte buffer, then parsed directly from
that buffer. The buffer is released before processing the next file. A single
application read request may still require multiple operating-system or device
transfers.
\resultplot{task-2}{Task 2: elapsed time, speed-up, and peak resident memory}{fig:2}
\begin{table}[H]\centering
\caption{Task 2 results; one measurement per configuration}
\begin{tabular}{rrrr}\hline
Threads & Time (s) & Speed-up & Peak RSS (MiB) \\\hline
1 & 11.174 & 1.000 & 47.97 \\
2 & 8.177 & 1.366 & 92.16 \\
4 & 6.522 & 1.713 & 180.22 \\\hline
\end{tabular}
\end{table}

\section{Comparison and limiting resources}
Task 2 on one thread takes 11.174 s, compared with 9.942 s for Task 1a on one
thread: a difference of 1.231 s, or about 12.4\%. Dividing total input bytes by
elapsed time gives approximately 33.1 MB/s and 37.2 MB/s respectively. These
are end-to-end processing rates, not isolated device read speeds.

The comparison changes both cache policy and reading strategy: Task 1a uses
warm-cache line streaming, while Task 2 reads whole files under the recorded
reconnect policy. It does not isolate USB latency. Moreover, the USB model and
link version are absent. The available numbers therefore support a mixed
storage-and-parsing workload, but do not conclusively establish that USB I/O
alone dominates execution time.

More threads do help in this series: two threads reduce time by about 26.8\%
and four by about 41.6\%, relative to one. Four threads achieve 1.713 times the
throughput, well below ideal fourfold scaling. Parsing different buffers can
proceed in parallel, and some computation can overlap reads. All reads still
share the same USB device and connection, so concurrency cannot multiply their
physical bandwidth. One sample per configuration limits conclusions about
repeatability.

\section{Memory cost}
Peak RSS rises approximately with the number of simultaneous buffers:
47.97 MiB, 92.16 MiB, and 180.22 MiB. An average split file occupies about
44.04 MiB, so four live file buffers explain most of the four-thread peak.
The implementation processes files sequentially within each worker, avoiding
retention of all eight buffers at once.

\chapter{Control questions}
\section{Physical cores and logical processors}
A physical core contains execution hardware; SMT exposes more than one logical
processor sharing that hardware. This workstation has four cores and eight
logical processors. Task 1a improves strongly up to four threads, then only
modestly from four to eight. This is consistent with shared core resources,
although the measurements alone do not isolate SMT's benefit.

\section{Correctness at file boundaries}
Internal byte offsets are advanced to complete-line boundaries. Adjacent
workers use the same boundary as an exclusive end and inclusive start.
Every line therefore lies in one range. The matching negative, zero, and
positive counters provide an empirical check across all recorded Task 1a
configurations, in addition to the boundary reasoning.

\section{Lost updates in shared counters}
Two threads can read the same old value, independently increment it, and write
the same new value. One increment is lost. The race experiment demonstrates
incorrect totals in all three runs, but its undefined behavior prevents
predicting exact results. A mutex makes each update mutually exclusive.

\section{Why local counters are faster}
Private counters avoid a lock on every element and repeated writes to the same
shared cache lines. Workers synchronize at completion, then a small reduction
adds three values per thread. This replaces synchronization proportional to the
50 million elements with a reduction proportional to the thread count.

\section{I/O-bound storage and network services}
An I/O-bound computation spends much of its elapsed time waiting for external
data rather than using available CPU capacity. A single USB drive has shared
bandwidth; additional readers eventually compete for the same resource.
A server handling independent network requests can often use otherwise idle
time to serve another request while the first waits for a remote response.
Concurrency helps when there is independent work and unused capacity; it is
not a guarantee for every I/O-bound system. Network bandwidth and remote
services can also saturate. Our USB measurements show partial benefit, so they
do not justify claiming that concurrency never helps USB access.

\section{Static partitions versus a queue}
Static partitioning assigns work before execution and avoids queue operations.
It works well when partitions have similar costs, as in this sign-counting
workload. A queue assigns the next batch to whichever consumer becomes free,
which can handle uneven task costs more flexibly. That flexibility requires
synchronization, batching, and memory for queued work. Here, independent
partitions achieve the best absolute time, while the queue adds substantial
memory and coordination costs.

\utmconclusions
The measurements show that workload structure matters more than the thread
count alone. Independent partitions reach 2.857 times their one-thread
throughput with little resident-memory growth. Most scaling occurs near the
physical-core count, but the recorded medians continue improving through
64 threads; a complete plateau is not observed for Task 1a.

Local counters are both correct and fastest in the array experiment. A shared
mutex preserves correctness at a large measured cost, while unsynchronized
updates produce incorrect counts in every run. The single-reader queue reaches
a practical plateau around two consumers and offers at most 1.092 times its
one-consumer throughput in the saved medians. This is consistent with sequential
production limiting consumer utilization. The additional six-reader experiment
has substantially lower recorded times, but unmatched measurement sessions
prevent attributing the entire difference to reader count alone.

Warm-cache file parsing and in-memory counting are consistent with CPU-side
work and synchronization costs. The USB series benefits from extra threads
but shows sublinear scaling and substantially greater memory use from whole-file
buffers. Its data do not isolate a purely I/O-bound bottleneck. Missing USB
identification, uncertain historical build flags, and the absence of resource
profiling limit stronger machine-specific conclusions.
\end{document}
